\subsection{The primitive Klein contribution inside the actual average polynomial}

\begin{proposition}[P-F44: an exact signed global cone contribution]
\label{prop:kleinglobalcontribution}
Let $n\ge4$ with $3\nmid n$, and retain $B_n$ and the exact
Weyl--Brion expansion of P-F42. Set
\[
 P=\frac{2^n-(-1)^n}{3},\qquad D=\binom n3,
\]
and let $C_a$ be the moving dual-normal determinant coefficient
of the Klein cubic in P-F40, for $a\in(\mathbb Z/P\mathbb Z)^\times$.
There is the following exact decomposition of the \emph{actual}
average constituent polynomial:
\[
 B_n(X)=\mathcal K_n+\mathcal R_n(X),\qquad
 \mathcal K_n=
   \frac1{nP}\sum_{\substack{a\bmod P\\(a,P)=1}}C_a
   =\frac1{nP}\operatorname{Tr}_{\mathbb Q(\zeta_P)/\mathbb Q}(C_1).
 \tag{P-F44a}
\]
Here $\mathcal R_n$ is given explicitly below by all the other
balanced character terms in the global Weyl--Brion expansion.
The removed terms are exactly the primitive Klein characters
on directed Hamiltonian-cycle bases, with all three central
cube-root lifts included.
They have degree zero as polynomials in $X$, and their total
is independent of the perturbation and specialization direction
used in P-F42.

If $n\equiv3\pmod4$, then
\[
                              \mathcal K_n=0.
 \tag{P-F44b}
\]
Thus the primitive Klein terms cancel in this particular
part of the actual global average polynomial, despite the
large individual amplitudes in P-F40.
For even $D$, formula (P-F44a) is a signed trace, and no
nonvanishing, sign, or dominance over $\mathcal R_n$ is claimed.
The normalization is $1/(nP)$, rather than the $1/P$ appearing
in the cyclic normal-model Molien sum.
\end{proposition}

\begin{proof}
We first identify a general collection of degree-zero cone terms.
Use the matrix $A$, the Weyl shifts $b_\sigma=\delta-\sigma\delta$,
and the character-polynomials $P_{J,\chi,b_\sigma}$ of P-F42.
Write $z_i=\chi(e_i)$ for a character of
$\mathbb Z^n/A_J\mathbb Z^n$.
Suppose the following hold:
\[
 \begin{gathered}
 A_J^{-1}\mathbf1\ \hbox{has all coordinates strictly positive},
 \qquad \chi(3\mathbf1)=1,\\
 z_1,\ldots,z_n\ \hbox{are distinct},\qquad
 \{\alpha:z^\alpha=1,\ |\alpha|=3\}=J.
 \end{gathered}
 \tag{P-F44c}
\]
The character is balanced because $c=3$ when $3\nmid n$.
The first condition means that $J$ is eventually feasible
for every Weyl shift and every fixed perturbation allowed
in P-F42. The last condition means that no nonbasic
denominator factor has character one, so the number $k$
in that proof is zero.
Consequently (P-F42l) has a constant term independent of
the stretch variable and the specialization direction:
\[
 P_{J,\chi,b_\sigma}(X)=
 \frac{z^{-b_\sigma}}
 {d_J\prod_{\alpha\notin J}(1-z^\alpha)},\qquad
 d_J=|\det A_J|.
 \tag{P-F44d}
\]
It is retained by averaging, since the character is balanced.

Put
\[
 \Delta(z)=z^{-\delta}
                  \sum_{\sigma\in S_n}
                         \operatorname{sgn}(\sigma)z^{\sigma\delta}
          =\prod_{i<j}(1-z_j/z_i).
 \tag{P-F44e}
\]
Summing (P-F44d) over the signed Weyl shifts therefore gives
\[
 T(J,z)=
 \frac{\Delta(z)}
 {d_J\prod_{\alpha\notin J}(1-z^\alpha)}.
 \tag{P-F44f}
\]
Permuting the variables permutes pairs $(J,z)$, preserves
all the conditions in (P-F44c), and preserves the denominator
product and $d_J$. This action on pairs is free because the
$z_i$ are distinct.
The elementary alternant identity
\[
 \sum_{\tau\in S_n}\Delta(\tau z)
       =\prod_{i\ne j}(1-z_i/z_j)
 \tag{P-F44g}
\]
then shows that a complete variable-permutation orbit of
such pairs contributes exactly
\[
 \frac1{d_J}\,
 \frac{\prod_{i\ne j}(1-z_i/z_j)}
      {\prod_{\alpha\notin J}(1-z^\alpha)}.
 \tag{P-F44h}
\]
For clarity, if
$A_\delta(z)=\sum_\sigma\operatorname{sgn}(\sigma)
z^{\sigma\delta}$, the sum in (P-F44g) is
$A_\delta(z)A_\delta(z^{-1})$.
The Vandermonde product gives precisely the right-hand
side of (P-F44g), including its sign.

We now apply this calculation to the Klein bases.
A directed Hamiltonian cycle on the $n$ labelled variables
gives the $n$ monomials $x_i^2x_{\nu(i)}$, where $\nu$ is
an $n$-cycle. There are $(n-1)!$ distinct bases of this
form. For the standard cycle, with indices $0,\ldots,n-1$,
the basis matrix is $2I$ plus a cyclic permutation matrix.
Thus
\[
 d_J=2^n-(-1)^n=3P,\qquad
                  A_J^{-1}\mathbf1=\tfrac13\mathbf1.
 \tag{P-F44i}
\]
The determinant formula follows directly from the two
nonzero terms in its determinant expansion.

Let $\rho=e^{2\pi i/P}$ and $w_i=(-2)^i$.
For each $\omega\in\mu_3$ and each unit $a\bmod P$, take
\[
                   z_i=\omega\rho^{a w_i}.
 \tag{P-F44j}
\]
These are $3\varphi(P)$ distinct characters for the
standard basis. Distinctness follows from
$\mu_3\cap\mu_P=\{1\}$, since $3\nmid P$.
Indeed $(-2)^n=(1-3)^n\equiv1-3n\pmod9$, so
$(1-(-2)^n)/3\equiv n\pmod3$; this integer is $P$
up to sign, and $3\nmid n$.
They annihilate every basis monomial and hence are
characters of the finite group of order $3P$.
Their determinant is $\omega^n$, because the Klein
matrix in P-F40 belongs to $SL_n$. In particular
$\chi(3\mathbf1)=1$.
P-F40 proves that their eigenvalues are distinct and
their fixed cubic monomials are exactly the displayed
$n$ cycle monomials. Multiplication by $\omega$ has no
effect on a cubic monomial character. Thus every
condition in (P-F44c) holds.
Transport these characters by variable permutations to
every directed-cycle basis. Denote the resulting set
of pairs by $\mathscr K_n$.

We count its variable-permutation orbits precisely.
The subgroup preserving the standard directed-cycle
basis is the cyclic group of $n$ rotations. It acts
on $a$ by multiplication by powers of $-2$.
The order of $-2$ modulo $P$ is exactly $n$: its
$n$th power is one, and the distinct Klein eigenvalues
proved in P-F40 exclude a smaller positive exponent.
Thus each orbit on units has size $n$.
The central index $\omega$ cannot be mixed by a variable
permutation: it is determined by the determinant
$\omega^n$, since $3\nmid n$.
It follows that the orbits of $\mathscr K_n$ under all
of $S_n$ are indexed by
\[
 \omega\in\mu_3,\qquad
 a\in(\mathbb Z/P\mathbb Z)^\times/\langle-2\rangle.
 \tag{P-F44k}
\]
Equivalently,
$|\mathscr K_n|=3(n-1)!\varphi(P)$ and every full
permutation orbit has size $n!$, giving the same count.

The determinant quotient in (P-F44h) is independent of
$\omega$. For the Klein cubic, exactness of the tangent
sequence identifies it with the reciprocal moving
\emph{primal} normal determinant. P-F40 defines $C_a$
using the dual normal representation, so the quotient
is $C_{-a}$, rather than $C_a$:
\[
 \frac{\prod_{i\ne j}(1-z_i/z_j)}
      {\prod_{\alpha\notin J}(1-z^\alpha)}
       =\frac1{\det(1-\beta^a\mid S_{\rm mov})}
       =C_{-a}.
 \tag{P-F44l}
\]
This explicitly accounts for dual conventions and
possible signs.
The quantity $C_a$ is constant on a class modulo
multiplication by $-2$, because rotation conjugates
the corresponding Klein automorphisms.
Combining (P-F44h)--(P-F44l), the total contribution
of $\mathscr K_n$ is
\[
 \sum_{\omega\in\mu_3}
 \sum_{[a]\bmod\langle-2\rangle}\frac{C_{-a}}{3P}
   =\frac1{nP}\sum_{(a,P)=1}C_{-a}
   =\frac1{nP}\sum_{(a,P)=1}C_a.
 \tag{P-F44m}
\]
The last equality reindexes units by negation.
Each $C_a$ is the Galois conjugate of $C_1$ under
$\rho\mapsto\rho^a$, which proves the trace formula
and rationality in (P-F44a).
P-F40 proves $C_{-a}=(-1)^D C_a$.
Since $D$ is odd exactly when $n\equiv3\pmod4$,
pairing units under negation proves (P-F44b).
Here $P>2$, so no unit is its own negative.

To specify the remainder without making any local-to-global
identification, choose any perturbation and common direction
as in P-F42. For a Weyl shift $\sigma$, let
$\mathscr B_\sigma$ be its eventually feasible bases.
Then
\[
 \mathcal R_n(X)=
 \sum_{\sigma\in S_n}\operatorname{sgn}(\sigma)
 \sum_{J\in\mathscr B_\sigma}
 \sum_{\substack{\chi\in\widehat{\mathbb Z^n/A_J\mathbb Z^n}\\
                 \chi(3\mathbf1)=1,\ (J,\chi)\notin\mathscr K_n}}
              P_{J,\chi,b_\sigma}(X).
 \tag{P-F44n}
\]
All sets and polynomial terms in this expression are
exactly those already constructed in P-F42.
Every Klein basis occurs in every $\mathscr B_\sigma$
by (P-F44i), so removing precisely $\mathscr K_n$
from that global expansion removes the constant
(P-F44m) and nothing else.
This proves the asserted decomposition of $B_n$.
Although the displayed representation of the remainder
uses the chosen perturbation and direction, its value
is the intrinsic polynomial $B_n-\mathcal K_n$.
\end{proof}

\paragraph{What this does and does not settle.}
Formula (P-F44a) supplies a global comparison for this
specified collection of terms, derived from the exact
lattice-point expansion. It does not identify the cyclic
normal-model Hilbert series with the global invariant
Hilbert series.
The large value of an individual $C_a$ does not bound
its Galois trace, and the trace can still cancel against
the constant term of $\mathcal R_n$.
More explicitly,
\[
 U_n=
   \bigl|\mathcal K_n+[X^0]\mathcal R_n\bigr|
       +\sum_{j\ge1}\bigl|[X^j]\mathcal R_n\bigr|.
 \tag{P-F44o}
\]
Thus the remaining quantitative gate is an explicitly
signed coefficient estimate for (P-F44n), together with
its possible constant-term cancellation with the trace.
No unsigned dominance or index lower bound is inferred.

\paragraph{A squarefree eigenbasis for the moving Klein normal module.}
There is a useful exact simplification of the coefficient in
(P-F44a). As a representation of the specified cyclic group
$\Gamma=\langle\beta\rangle$,
\[
                  S_{\rm mov}\cong\bigwedge\nolimits^3\mathbb C^n,
 \qquad
 C_a=
 \prod_{0\le i<j<k<n}
       \left(1-\rho^{-a(w_i+w_j+w_k)}\right)^{-1}.
 \tag{P-F44p}
\]
The exterior power here carries the diagonal action with
eigenvalues $\rho^{w_i}$ on $\mathbb C^n$.
The isomorphism is a statement about this cyclic stabilizer
representation, not an assertion of an $SL_n$-module structure
on a normal space at a nonfixed point.

Here is a direct proof, including the cyclic boundary cases.
Let $T=\mathbb C[x_0,\ldots,x_{n-1}]$ and let $J_K$ be the
Jacobian ideal of $K_n$. Its generators are
\[
                 x_i^2+2x_{i+1}x_{i+2},
                         \qquad i\bmod n.
 \tag{P-F44q}
\]
The smoothness proof in P-F40 shows that the partial
derivatives have common zero only at the origin.
They form a regular sequence of quadrics. There are
therefore no relations with linear coefficients among
them: equivalently, the degree-three part of the first
Koszul syzygy module is zero, as each basic Koszul
syzygy already has degree four. Hence
\[
                    \dim(T/J_K)_3
                       =\binom{n+2}{3}-n^2=D.
 \tag{P-F44r}
\]

Every degree-three monomial reduces modulo $J_K$ to a
linear combination of squarefree monomials, in fact to
a scalar multiple of one or to zero. Cubes satisfy
$x_i^3=-2x_ix_{i+1}x_{i+2}$.
For $x_i^2x_j$ with $j\ne i$, use
$x_i^2x_j=-2x_{i+1}x_{i+2}x_j$.
If $j$ is different from $i+1,i+2$, this is squarefree.
When $j=i+1$, put $u_i=x_i^2x_{i+1}$; then
$u_i=-2u_{i+1}$, and going once around the cycle gives
$u_i=(-2)^n u_i$. Thus $u_i=0$ in characteristic zero.
When $j=i+2$, a second reduction gives
\[
                  x_i^2x_{i+2}
                     =4x_{i+1}x_{i+3}x_{i+4}.
 \tag{P-F44s}
\]
These three indices are distinct modulo every $n\ge4$:
their pairwise positive differences are one, two, and
three. In particular, at $n=4$ the final index is $i$,
and the three remaining indices are still distinct.
There are exactly $D$ squarefree cubic monomials, so
the spanning argument and (P-F44r) make their classes
a basis.

The degree-three Jacobian subspace is the image
$\mathfrak{gl}_n\cdot K_n$.
The absence of linear syzygies makes this Lie-orbit
map injective. Its subspace
$\mathfrak{sl}_n\cdot K_n$ has codimension one, with
the complementary invariant Euler line represented
by $K_n$ itself.
It follows that the quotient normal module $S$ in
P-F40 is an extension of $(T/J_K)_3$ by a trivial
one-dimensional representation.
Representations of the finite group $\Gamma$ over
$\mathbb C$ split by averaging.
Moreover P-F40 proves $\dim S^\Gamma=1$, so
$(T/J_K)_3$ has no fixed vectors and is exactly the
moving module.
Each squarefree basis vector has eigenvalue
$\rho^{w_i+w_j+w_k}$. Sending
$e_i\wedge e_j\wedge e_k$ to its squarefree class is
therefore a $\Gamma$-equivariant vector-space
isomorphism. This proves (P-F44p), including the
negative exponent required for the dual coefficient.

Thus the specified global constant contribution is
the completely explicit cyclotomic trace
\[
 \mathcal K_n=\frac1{nP}
 \operatorname{Tr}_{\mathbb Q(\rho)/\mathbb Q}
 \left[
   \prod_{0\le i<j<k<n}
          (1-\rho^{-(w_i+w_j+w_k)})^{-1}
 \right].
 \tag{P-F44t}
\]
This expression does not remove the signed trace or the
remaining polynomial in (P-F44n), and it is not used to
claim a lower bound for $U_n$.

\paragraph{Three exact finite trace values, including a negative contribution.}
The specified global contributions in three small cases are
\[
              \mathcal K_4=\frac1{25},\qquad
              \mathcal K_5=\frac2{121},\qquad
              \mathcal K_8=-\frac{624}{85}.
 \tag{P-F44u}
\]
For $n=4$, $P=5$ and the four dual squarefree weights
are all the nonzero residues modulo five.
For $n=5$, $P=11$ and the variable weights are
$1,9,4,3,5$ modulo eleven. Their sum is zero, and their
ten pair sums are the ten nonzero residues. Thus the
ten dual squarefree triple weights are also all the
nonzero residues. In both cases
$C_a=1/P$, since
$\prod_{j=1}^{P-1}(1-\rho^j)=P$.
Equation (P-F44a) gives the first two displayed values.

Here is a complete finite polynomial certificate for the
third value. Put
\[
 \Phi(x)=\Phi_{85}(x)
       =\frac{(x^{85}-1)(x-1)}{(x^5-1)(x^{17}-1)},\qquad
 D_8(x)=\prod_{0\le i<j<k<8}(1-x^{r_{ijk}}),
\]
where $r_{ijk}$ is the representative in $\{1,\ldots,84\}$
of $-((-2)^i+(-2)^j+(-2)^k)$ modulo $85$.
Define $N(x)=\sum_{j=0}^{63}N_jx^j$ by the following
table; a row labelled $j$ gives $N_j,\ldots,N_{j+7}$:
\[
\begin{array}{r|rrrrrrrr}
 j&N_j&N_{j+1}&N_{j+2}&N_{j+3}&N_{j+4}&N_{j+5}&N_{j+6}&N_{j+7}\\\hline
 0&-1167&0&1048&-808&0&708&0&454\\
 8&0&-708&254&0&0&-708&100&-340\\
16&0&708&-708&340&-100&0&708&100\\
24&-354&-340&-708&708&100&100&-694&100\\
32&708&-708&354&-694&340&-100&-708&808\\
40&354&-354&708&-708&-100&354&100&0\\
48&-708&0&-694&354&0&340&0&-694\\
56&808&808&-254&0&-1048&-100&808&-708
\end{array}
\]
Exact integer polynomial multiplication and monic division give
\[
                         D_8(x)N(x)\equiv17\pmod{\Phi(x)}.
 \tag{P-F44v}
\]
For a reproducible check requiring no symbolic inversion,
first obtain $\Phi$ by dividing
$1+x^5+\cdots+x^{80}$ by $1+x+\cdots+x^{16}$.
Starting at one, multiply the $56$ factors defining $D_8$,
reducing by this monic polynomial after each multiplication.
Multiply the result by the displayed $N$ and reduce once
more; the remainder is the constant $17$.
The independent integer-arithmetic verifier
\texttt{verify\_klein85\_certificate.py} implements exactly
these divisions and was executed successfully.
All the inputs to this finite identity are displayed above.
It follows that $C_1=N(\rho)/17$ for a primitive
$85$th root $\rho$.

The trace of $\rho^j$ is $64$ for $j=0$, is $-4$ for a
positive multiple of five less than $64$, is $-16$ for
a positive multiple of seventeen less than $64$, and
is one for the remaining positive indices below $64$.
Indeed, summing over residues coprime to $85$ and using
inclusion--exclusion on multiples of five and seventeen
gives
\[
 \sum_{(a,85)=1}\rho^{aj}
   =85\,\mathbf1_{85\mid j}
          -17\,\mathbf1_{17\mid j}
          -5\,\mathbf1_{5\mid j}+1.
\]
The table gives respectively
$N_0=-1167$, the three other coefficient-group sums
$-2934$, $1416$, and $744$. Hence
\[
 \operatorname{Tr}(C_1)
  =\frac{64(-1167)-4(-2934)-16(1416)+744}{17}
  =-4992.
 \tag{P-F44w}
\]
Dividing by $nP=8\cdot85$ proves the last value in
(P-F44u).
It is an exact negative contribution to the specified
part of the actual average polynomial.
It is not asserted to equal its full constant coefficient,
and gives no value or bound for $U_8$.
